\section{Compilation problem and scope}
\label{sec:introduction}
A conjunctive query specifies satisfying valuations, not a relation container,
index layout, or access algorithm. MiniLinq constructs a relational plan,
records required index demands, and then lets a caller-supplied,
query-independent policy fix one physical operator per result and bind every
keyed read to a host key fact or a Rust structure built over the result. A
printer emits one pull program over the runtime operation traits; the compiler
uses open operation traits. Supplied structure selection chooses concrete default
Rust types when the caller supplies no applicable type advice.

The compiled chain is
\[
 \mathsf{CQ}\longrightarrow\mathsf{RA}\longrightarrow\mathsf{IX}
 \xrightarrow{\;\text{policy}\;}\mathsf{Physical}
 \xrightarrow{\;\text{policy}\;}\mathsf{Structures}
 \longrightarrow\mathsf{Rust}.
\]
IX preserves the RA plan and adds named index demands.
Physical derives demand uses from its operands and restates each result with one
physical operator and the mode its reads allow. Lowering erases Physical choices
to RA and separately checks IX demands; Structures extends Physical unchanged
and supplies storage bindings. Section~\ref{sec:planning} expands the logical
arrow into normalization, ordering, optional reduction and elaboration.
Section~\ref{sec:pass-contracts} specifies each language's local contract,
including Physical and Structures, and the separate pass correspondence checks.

The retained relational theory covers valuations, named RA, host-expression
dependencies, GYO and Yannakakis planning. These logical results do not choose
physical storage. Rust quotations preserve host syntax; moving or duplicating
host computations additionally requires the totality, purity and binding
assumptions stated below. Recursive programs, aggregation, negation, general
physical cost search remain outside the supported contracts. Pull realization
and integrated generated-query execution are implemented.
